\section{Parallel Algorithms}

\subsection{Parallel Patterns}

\begin{definition}
A \textbf{reduction} produces a single answer from a collection via an associative operator.
\end{definition}

\begin{example}
Max, Min, Sum, Product of an array.
\end{example}

\begin{definition}
A \textbf{map} applies a function to each elements of a collection producing an output of the same size.
\end{definition}

\begin{example}
Vector addition: $v_i = u_i + w_i \forall i \in [1, n]$
\end{example}

\begin{definition}
A \textbf{stencil} extends the map pattern by using the element itself and its neighbours.  
\end{definition}

\begin{example}
A smoothing kernel (ex. 3x3) that calculates the average value is sled over an image to smooth/blur the image.
\end{example}

\subsection{Parallel Algorithms}

\subsubsection{Prefix-Sum}

Prefix-Sum: Sum of first $n$ inputs.

\begin{algorithm}
upTask(); //divide problem and build tree with sums of subproblems
downTask(); //fill in the correction values (left_sum + parent_correction)
\end{algorithm}

\begin{theorem}
The span of the prefix sum optimization is $O(\log n)$, therefore parallelism of $\frac{n}{\log n}$.
\end{theorem}

\subsubsection{Pack}

\textbf{Pack} pack all elements from an array that satisfy a certain condition into a seperate output. (by using prefix sum)

\begin{algorithm}
applyMap(); //applies a map to the input
//results in a binary vector [1, 0, 1, 1, 0, ...]\\
parallelPrefixSum(); //calculates the prefix sum of the binary vector\\
//results in [1, 1, 2, 3, 3, ...]\\
applyPackMap(); //applies another map to the result of the prefix sum\\
//results in [0, 2, 3, ...] (indices of elements to write to output)
\end{algorithm}

\subsubsection{Quicksort}

We can reduce the \textbf{limiting factor of the problem division} from a span of $O(n)$ to a span of $O(\log n)$
using the pack-pattern.

\begin{algorithm}
applyMaps(); //apply maps for greater and smaller than pivot
parallelPrefixSum(); //count how many smallerthan/greatherthan elements come bofore
applyMap(); //index = numLessThan || numTotalLessThan + numGreaterThan
\end{algorithm}

We can do this by using two pack-patterns to partition the data. This leads to a partitioning of $O(\log n)$
levels times $O(\log n)$ for prefixSum hence $O(\log^2 n)$ total span. 

This means a parallelism of $\frac{n\log n}{\log^2 n} = \frac{n}{\log n}$.

        